\scanpage{61}

\begin{center}- 152 -\end{center}
\begin{flushright}IX\end{flushright}

$\uline{\operatorname{Pic}}^{\unsure{o|0}}_{\widetilde C/k}$ est de ``rang réductif'' nul (i.e. son tore maximal est nul), de sorte que le tore maximal de $\uline{\operatorname{Pic}}^{\unsure{o|0}}_{C/k}$ est contenu dans
\[
N=\operatorname{Ker}(\uline{\operatorname{Pic}}_{C/k}
\longrightarrow\uline{\operatorname{Pic}}_{\widetilde C/k})\quad,
\]
donc est isomorphe canoniquement au tore maximal de ce dernier groupe. Or la structure de $N$ s'explicite aisément à l'aide de EGA IV 21.8.5 : on trouve que $N$ est canoniquement une extension d'un tore $T$ dont on va expliciter la description, par un groupe unipotent connexe $U$ (qui se dévisse de même, sur $k$, avec une suite de composition à quotients isomorphes au groupe additif $G_a$). D'ailleurs, si $T'$ est le tore maximal de $N$, alors $T'\cap U$ est à la fois unipotent et de type multiplicatif, donc nul (SGA 3 XVII 2.4), et l'image de $T'$ dans $T$ est identique à $T$ (SGA 3 XII 7.1 e)), de sorte que $T'\longrightarrow T$ est un isomorphisme. Donc $T$ est canoniquement isomorphe au tore maximal de $N$, donc de $\uline{\operatorname{Pic}}^{\unsure{o|0}}_{X/k}$.

C'est pour expliciter la structure de $N$ qu'il sera commode de supposer $k$ séparablement clos, de sorte que les composantes irréductibles $C_i$ de $C$ sont géométriquement irréductibles, et que les points non normaux de $C$ sont radiciels sur $k$. Ceci dit, soit $J$ l'ensemble des composantes irréductibles de $\widetilde C$, et pour tout $x\in C$, soit $J(x)$ l'ensemble des ``branches'' de $C$ passant par $x$ (i.e. des points de $\widetilde C$ au-dessus de $x$); soient enfin
\[
(12.3.1)\qquad R(x)=\mathbb Z^{J(x)}/\operatorname{Im}\mathbb Z\quad,
\]
où $\mathbb Z\longrightarrow\mathbb Z^{J(x)}$ est l'application diagonale, et
\[
(12.3.2)\qquad R=(\bigoplus_x R_x)/\operatorname{Im}\mathbb Z^J\quad,
\]

\margin{470}
\margin{470}
\margin{60}
